% This LaTeX document needs to be compiled with XeLaTeX.
\documentclass[
  10
]{article}
\usepackage[margin=2cm,includehead=true,includefoot=true,centering,]{geometry}
\usepackage{xcolor}
\usepackage{ucharclasses}
\usepackage{hyperref}
\usepackage{amsmath,amssymb}
\usepackage{amsfonts}
\usepackage[version=4]{mhchem}
\usepackage{stmaryrd}
\usepackage{bbold}
\usepackage{polyglossia}
\usepackage{fontspec}
\usepackage{ctex}
\usepackage[export]{adjustbox}
\usepackage{tabularx}
\usepackage{booktabs}

\usepackage{setspace}
\setstretch{1.2}

\makeatletter
\@ifundefined{KOMAClassName}{% if non-KOMA class
  \IfFileExists{parskip.sty}{%
    \usepackage{parskip}
  }{% else
    \setlength{\parindent}{0pt}
    \setlength{\parskip}{6pt plus 2pt minus 1pt}}
}{% if KOMA class
  \KOMAoptions{parskip=half}}
\makeatother
\usepackage{longtable,booktabs,array}
\newcounter{none} % for unnumbered tables
\usepackage{multirow}
\usepackage{calc} % for calculating minipage widths
% Correct order of tables after \paragraph or \subparagraph
\usepackage{etoolbox}
\makeatletter
\patchcmd\longtable{\par}{\if@noskipsec\mbox{}\fi\par}{}{}
\makeatother
% Allow footnotes in longtable head/foot
\IfFileExists{footnotehyper.sty}{\usepackage{footnotehyper}}{\usepackage{footnote}}
\makesavenoteenv{longtable}
\usepackage{graphicx}
\makeatletter
\newsavebox\pandoc@box
\newcommand*\pandocbounded[1]{% scales image to fit in text height/width
  \sbox\pandoc@box{#1}%
  \Gscale@div\@tempa{\textheight}{\dimexpr\ht\pandoc@box+\dp\pandoc@box\relax}%
  \Gscale@div\@tempb{\linewidth}{\wd\pandoc@box}%
  \ifdim\@tempb\p@<\@tempa\p@\let\@tempa\@tempb\fi% select the smaller of both
  \ifdim\@tempa\p@<\p@\scalebox{\@tempa}{\usebox\pandoc@box}%
  \else\usebox{\pandoc@box}%
  \fi%
}
% Set default figure placement to htbp
\def\fps@figure{htbp}
\makeatother
\setlength{\emergencystretch}{3em} % prevent overfull lines
\providecommand{\tightlist}{%
  \setlength{\itemsep}{0pt}\setlength{\parskip}{0pt}}


\hypersetup{colorlinks=true, linkcolor=blue, filecolor=magenta, urlcolor=cyan,}
\urlstyle{same}

\usepackage{colortbl}
\definecolor{table-row-color}{HTML}{999999}
\definecolor{table-rule-color}{HTML}{999999}
%\arrayrulecolor{black!40}
\arrayrulecolor{table-rule-color}     % color of \toprule, \midrule, \bottomrule
\setlength{\aboverulesep}{0pt}
\setlength{\belowrulesep}{0pt}


\setotherlanguages{english}
\IfFontExistsTF{Source Han Serif CN}
{\newfontfamily\chinesefont{Source Han Serif CN}}
{\IfFontExistsTF{Noto Serif CJK SC}
  {\newfontfamily\chinesefont{Noto Serif CJK SC}}
  {\IfFontExistsTF{SimSun}
    {\newfontfamily\chinesefont{SimSun}}
    {\IfFontExistsTF{FangSong}
      {\newfontfamily\chinesefont{FangSong}}
      {\newfontfamily\chinesefont{Arial Unicode MS}}
}}}
\IfFontExistsTF{Times New Roman}
{\newfontfamily\englishfont{Times New Roman}}
{\IfFontExistsTF{Liberation Serif}
  {\newfontfamily\englishfont{Liberation Serif}}
  {\IfFontExistsTF{DejaVu Serif}
    {\newfontfamily\englishfont{DejaVu Serif}}
    {\newfontfamily\englishfont{Arial}}
}}


\author{}
\date{}

\begin{document}


\section{Enhancing Spatiotemporal Consistency in Intracranial Aneurysm
Detection: ATargeted Feature Aggregation Framework for DSA
Sequences}\label{enhancing-spatiotemporal-consistency-in-intracranial-aneurysm-detection-atargeted-feature-aggregation-framework-for-dsa-sequences}

Anonymous submission

\section{Abstract}\label{abstract}

Intracranial aneurysm (IA) rupture is the leading cause ofnon-traumatic
subarachnoid hemorrhage, resulting in severedisability or death. Digital
subtraction angiography (DSA)is the gold standard for diagnosing
cerebrovascular diseasesand plays a vital role in early IA detection.
However, accu-rate IA identification in DSA sequences presents two
pri-mary challenges: small IA signals are obscured by redun-dant
background features, and detection results frequentlylack consistency
across the temporal sequence. To addressthese challenges, a detection
framework with targeted fea-ture aggregation is proposed. This approach
refines featuresby selectively processing information only from
high-qualityproposals. This targeted mechanism not only effectively
sup-presses background noise but also leads to more
consistentdetections, enabling comprehensive temporal analysis.
Fur-thermore, we employ a post-processing heuristic named Bi-SeqNMS,
which utilizes bidirectional temporal linking to fur-ther optimize
detection consistency across the sequence. Ex-perimental results show
our method achieves a significant im-provement over baseline models. Our
method demonstratesrobust performance on IAs with varied
characteristics, con-firming its potential in assisting IA diagnosis.

\section{1 Introduction}\label{introduction}

Intracranial aneurysm (IA), a cerebrovascular disease withsignificant
risks of disability and mortality, has a reportedprevalence of \(3 \%\)
to \(7 \%\) in recent years (Li et al.~2013). IArupture is the leading
cause of non-traumatic subarachnoidhemorrhage, accounting for most
disability and mortalitycases (Ropper and Zervas 1984). Given these
severe conse-quences, early and accurate detection of IAs combined
withtimely interventions is essential to mitigate patient risks.Magnetic
resonance angiography (MRA), computed tomog-raphy angiography (CTA), and
digital subtraction angiogra-phy (DSA) are commonly employed techniques
for diagnos-ing IAs in clinical practice (Hiratsuka et al.~2008; Yoon et
al.2016). For preoperative evaluation, clinicians primarily relyon MRA
and CTA images for the rapid screening of intracra-nial vascular lesions
(Piotin et al.~2003). Compared to CTAand MRA, DSA provides clearer,
higher-resolution vascu-lar imaging. Consequently, when CTA and MRA lack
suf-ficient diagnostic accuracy, the superior resolution of DSAenables
precise identification and localization of vascular le-sions,
establishing it as the international gold standard for di-

\pandocbounded{\includegraphics[keepaspectratio,alt={image}]{images/466d1a9cd3897ba07e2f20a2e6efc7d73eb001f79e4df0ee27cceef8cc75ce4e.jpg}}

\pandocbounded{\includegraphics[keepaspectratio,alt={image}]{images/885b4b79fc0eef33cedda1a9e4a21f0bb460414b87550e9e6108fa81d8da83bd.jpg}}

\pandocbounded{\includegraphics[keepaspectratio,alt={image}]{images/1e592f4d9ba1dc1b6779675f748bcf70a347928ad543ae1f7301fc72192e721e.jpg}}

\pandocbounded{\includegraphics[keepaspectratio,alt={image}]{images/5275891ae38c81f93cd7d0385e315e96c5745378c562b6fdccebb7efb46adad5.jpg}}

\pandocbounded{\includegraphics[keepaspectratio,alt={image}]{images/e54f1a57dc58f8e686071a096fa72ec2937c7cd0ac06a516e44b063bb3765dff.jpg}}

\pandocbounded{\includegraphics[keepaspectratio,alt={image}]{images/b2c1d9a20a806e27242d55cc44b35d879e3c3f838a86ee2a186366656977b00b.jpg}}

Figure 1: DSA sequence segment containing 6 frames fromthe arterial
phase. The red arrow indicates confirmed IA.This segment illustrates the
IA's progression from being ob-scured by surrounding vasculature to its
eventual emergence.

agnosing cerebrovascular diseases (Rodriguez-Regent et al.~´2014).
Despite its diagnostic superiority, DSA requires clin-icians to analyze
sequential images frame-by-frame. Thistime-consuming and labor-intensive
process imposes sub-stantial burdens on clinicians. Prolonged review may
exac-erbate subjective biases among clinicians, potentially
com-promising diagnostic accuracy (Shen, Wu, and Suk 2017).

Computer-Aided Diagnosis (CAD) systems effectively al-leviate
clinicians' workload and reduce diagnostic variabil-ity by providing
preliminary diagnostic results. Particularlyin aneurysm diagnosis, deep
learning-based methods havedemonstrated superior robustness in lesion
detection com-pared to conventional image processing techniques,
estab-lishing them as a focal point of research. Shi et
al.~(2020)proposed an end-to-end 3D-CNN segmentation model basedon CTA
images to assist in diagnosing IAs. Faron et al.(2020) employed a
CNN-based DeepMedic framework todetect IAs from clinical TOF-MRA data,
achieving sensitiv-ity comparable to that of radiologists. These results
highlightthe remarkable success of CNN-based object detection mod-els in
the CAD field. However, current research on CAD sys-tems for lesion
detection based on DSA images remains rel-atively limited, likely due to
the challenges associated with

acquiring high-quality DSA lesion detection datasets.

Existing DSA lesion detection methods generally adopttwo paradigms:
single-frame analysis, which relies on im-ages with explicit lesion
annotations, and multi-frame anal-ysis that aggregates spatiotemporal
features across frames.For instance, Xu et al.~(2023) designed a
plug-and-playloss function for YOLOX that considers image similarity
be-tween predicted and ground-truth bounding boxes, therebyimproving
vascular stenosis detection performance on spe-cific DSA frames. Mo et
al.~(2023) proposed a graph convo-lution module applied to Mask2Former
to dynamically re-assign label information and adaptively adjust
confidencefor erroneous instances, enabling automated segmentationof IAs
in DSA images. While these frame-level approacheshave demonstrated
notable potential, they are subject to in-herent limitations. Due to
factors such as variations in imag-ing angles and contrast agent flow,
intracranial arterial le-sions exhibit significant variability within
individual DSAsequences. These lesions may appear in an uncertain
num-ber of frames with diverse morphological features (Zhanget
al.~2025), as illustrated in Fig. 1. Consequently, relyingsolely on
individual DSA frames for inference cannot fullycapture the entire
cerebrovascular structure or these patho-logical characteristics.

Temporal feature learning architectures developed forvideo analysis
(Shi, Wang, and Guo 2023) have demon-strated their transferability to
medical image sequences.Duan et al.~(2019) first proposed a two-stage
CNN architec-ture to automatically detect IAs in specific vascular
regionsof 2D digital subtraction angiography (2D-DSA) sequences.Zeng et
al.~(2019) employed a spatial information fusionmethod for aneurysm
detection in 3D rotational angiography(3D-RA), in which consecutive
frames from each imagingtime series were extracted and concatenated to
embed tem-poral features into a single 2D image. Su et al.~(2022)
intro-duced a spatiotemporal deep learning framework for auto-mated
detection of intracranial vascular perforations in DSAsequences. This
framework utilized existing backbone net-works to extract features from
DSA frames and employeda long short-term memory (LSTM) network to
capture tem-poral dependencies between frame features. The study
con-firmed that aggregating more frames improves performance.However,
the substantial GPU memory demands associatedwith processing
high-resolution DSA images constitute abottleneck that constrains the
number of aggregable frames.To address this issue, study (Zhang et
al.~2023) employedsemi-automatic region-of-interest (ROI) extraction to
reducethe input frame size for the model, thereby allowing moreframes to
be aggregated. However, this approach still neces-sitated human
intervention during inference.

While Conv-LSTM architectures have proven effective incapturing temporal
dynamics from DSA sequences, theirreliance on processing complete
feature maps presents in-herent limitations for intracranial aneurysm
detection. Thechallenge arises from the spatial disparity between
lesionsize and image size. Clinical data as Fig. 1 indicates thatmost
IAs occupy less than \(5 \%\) of the frame area, result-ing in
aneurysm-specific features being diluted within exten-sive background
features. These dominant background fea-

Table 1: Statistical characteristics of our dataset. SPA, CPAand RA
stand for Sagittal Plane Angiograph, Coronal PlaneAngiograph and
Rotational Angiograph, respectively.

{\def\LTcaptype{none} % do not increment counter
\begin{longtable}[]{@{}llll@{}}
\toprule\noalign{}
\endhead
\bottomrule\noalign{}
\endlastfoot
Artery & Total & View & Diseases \\
\multirow{3}{*}{ICA} & \multirow{3}{*}{55} & SPA: 25 & Aneurysm: 39 \\
& & CPA: 20 & Stenosis: 16 \\
& & RA: 10 & \\
\multirow{3}{*}{ECA} & \multirow{3}{*}{43} & SPA: 17 & Aneurysm: 31 \\
& & CPA: 14 & Stenosis: 12 \\
& & RA: 12 & \\
\multirow{2}{*}{VA} & \multirow{2}{*}{30} & SPA: 15 & Aneurysm: 8 \\
& & CPA: 15 & Stenosis: 22 \\
\end{longtable}
}

tures not only obscure the discriminative signatures of le-sions (Zhang
et al.~2023), but also impose significant com-putational costs that
limit the number of frames available fortemporal aggregation.

Therefore, in this article, we propose a deep learningframework for
detection of IAs in two-dimensional digi-tal subtraction angiography
(2D-DSA) and unreconstructedthree-dimensional rotational angiography
(3D-RA) images.This framework is based on a carefully prepared and
anno-tated clinical DSA dataset, with its details presented in Ta-ble 1.
To achieve this, we developed a targeted feature ag-gregation approach
and a customized sequence-level post-processing method, while
extensively evaluating several ex-isting advanced techniques. The main
contributions of thispaper are summarized as follows:

• We propose a detection framework centered on a targetedfeature
aggregation approach, which integrates a featureselection process with a
self-attention aggregation mod-ule. By concentrating on high-quality
proposals, it mini-mizes background interference, thereby significantly
en-hancing detection accuracy and consistency.

• The framework is further enhanced by a bidirectionaltemporal linking
post-processing heuristic, referred to asBi-SeqNMS, to provide a final
refinement to the consis-tency established by targeted feature
aggregation.

• Comprehensive experiments conducted on our clinicalDSA dataset
demonstrate that the proposed frameworkachieves superior performance,
confirming its effective-ness as an accurate and reliable solution for
IA detection.

\section{2 Proposed method}\label{proposed-method}

In this section, we propose a framework for the detectionof IAs in DSA
sequences, as illustrated in Fig. 2. Theframework first employs a
baseline detector to generate ini-tial proposals. Subsequently, the
framework utilizes a self-attention-based Feature Aggregation Block
(FAB) to selec-tively aggregate spatiotemporal features from
high-qualityproposals, which enhances the target signal while
suppress-ing noise. Finally, a post-processing module that lever-ages
spatiotemporal consistency is applied to refine the re-

\pandocbounded{\includegraphics[keepaspectratio,alt={image}]{images/914d1db2474ffe80d09f630332bde442d3ecc7e36baf5a38ea0e5d577707e5b3.jpg}}

Figure 2: Proposed framework for IA detection, consisting of a base
detection model, feature selection process, feature ag-gregation block
and post-processing module named Bi-SeqNMS. The framework utilizes
enhanced DSA sequence images asinput and extracts features using the
base detection model. Feature selection process selects multiple
proposals along with theircorresponding features based on the detector's
scores for the predicted instances and the FAB optimizes these proposal
scores.In the prediction stage, the detection results from the FAB
undergo further temporal refinement through Bi-SeqNMS.

sults. Our targeted feature aggregation strategy and post-processing
module will be elaborated upon subsequently.

\section{2.1 Targeted Feature
Aggregation}\label{targeted-feature-aggregation}

In our proposed Targeted Feature Aggregation (TFA) ap-proach, we employ
a self-attention-based block that facili-tates learning of relationships
between features across dif-ferent DSA frames in the input. During IA
diagnosis inclinical DSA images, clinicians first identify specific
framesthat exhibit prominent IA features as evidence of aneurysmpresence
throughout the DSA sequence. Subsequently, forframes with less apparent
aneurysm characteristics, clini-cians determine IA locations based on
the temporal conti-nuity of the DSA sequence. This workflow aligns well
withthe self-attention mechanism (Vaswani et al.~2017), whichis a core
component of the Transformer architecture. Com-pared to RNNs and their
variants, this mechanism demon-strates superior long-range dependency
modeling capabilityfor capturing cross-frame relationships.

Feature Selection In previous studies on DSA lesion de-tection (Duan et
al.~2019; Su et al.~2022; Zhang et al.~2023),researchers typically
employed multi-scale feature maps ex-tracted by backbone networks for
multi-frame feature ag-gregation, leveraging inter-frame comparisons to
identifylesions. However, directly feeding complete feature maps

into FAB presents two significant limitations: First, exces-sive
background features may obscure the critical charac-teristics of IAs,
introducing interference in detection tasks;Second, the substantial
computational overhead associatedwith processing complete feature maps
severely restricts thenumber of frames that can be processed. To address
thesechallenges, we introduce a feature selection process prior
totestinputting the data into the FAB. By leveraging
YOLOv8'sTask-Aligned Learning (TAL) (Feng et al.~2021), we
rank\ldots predictions during training base on the alignment metric:

\[
t = s ^ {\alpha} \cdot u ^ {\beta}, \tag {1}
\]

where s represents the predicted class confidencescore, u denotes the
Complete Intersection over Union(CIoU) (Zheng et al.~2021) between
predicted andground-truth bounding boxes, and α, \(\beta\) are
balancinghyperparameters set to 0.5 and 6 respectively. Predictionsare
ranked by descending t, with the top-K candidatesselected as positive
samples. During the training phase,this metric guides our feature
screening. Specifically,predictions with t-values exceeding a threshold
are firstretained, followed by non-maximum suppression (NMS)to eliminate
redundant spatially overlapping proposals.Collected from the
classification branch, the selected ROIsand their corresponding features
\(f _ { i }\) ..are retained as inputsfor the FAB. During the inference
phase, since the ground-truth-dependent alignment metric t is
unavailable, we utilize

the classification confidence score \(s _ { i }\) as the quality
metricto rank proposals and select the top-K inputs for the
featureaggregation block.

Feature Aggregation As shown in Fig. 2, the input fea-tures of FAB are
defined as
\(F _ { i n } = \{ f _ { 1 } , { \bar { f } } _ { 2 } , \dots , f _ { n } \}\)
, whereeach \(f _ { i } \in \mathbb { R } ^ { n _ { i } \times d }\)
represents the feature set of the \(i\) -th frame.Here, \(n _ { i }\)
denotes the number of ROIs retained after thresholdfiltering and NMS for
the \(i\) -th frame, and \(d\) is the dimension-ality of ROI features.
All features in \(F _ { i n }\) are concatenatedinto a stacked matrix:

\[
F \in \mathbb {R} ^ {n \times d}, \quad \text {w h e r e} n = \sum_ {i = 1} ^ {n} n _ {i}. \tag {2}
\]

Through linear transformations, \(F\) is projected into query\(( Q )\) ,
key \(( K )\) , and value \(( V )\) matrices:

\[
Q = F W ^ {Q}, \quad K = F W ^ {K}, \quad V = F W ^ {V}, \tag {3}
\]

where
\(W ^ { Q } , W ^ { K } , W ^ { V } \in \mathbb { R } ^ { d \times d _ { k } }\)
. To enhance the model'sexpressive power, FAB employs multi-head
self-attention.The multi-head self-attention is computed as:

\[
h e a d _ {h} = \operatorname {s o f t m a x} \left(\frac {Q _ {h} K _ {h} ^ {\top}}{\sqrt {d _ {k}}}\right) V _ {h}, \tag {4}
\]

\[
\operatorname {M S A} (F) = \operatorname {C o n c a t} \left(\text {h e a d} _ {1}, \dots , \text {h e a d} _ {h}\right), \tag {5}
\]

with \(Q _ { h } , K _ { h } , V _ { h }\) denoting split sub-matrices
of the original\(Q , K , V\) . To preserve original feature information,
the ag-gregated features are concatenated with the original features:

\[
\operatorname {F A} (F) = \operatorname {C o n c a t} (\operatorname {M S A} (F), F), \tag {6}
\]

The aggregated feature vectors are subsequently processedthrough a
linear transformation followed a Softmax layer,yielding the refined
scores of FAB-optimized proposals.

During feature aggregation, it is essential to associateidentical IA
targets across frames. Drawing inspiration fromcosine similarity-based
cross-frame ROI association in two-stage video object detectors (Wu et
al.~2019), we proposematching proposal features using cosine similarity.
Specifi-cally, for the value features
\(V \in \mathbb { R } ^ { n \times d _ { k } }\) derived from
con-catenated \(F _ { i n }\) and linear projection, we compute the
row-normalized cosine similarity matrix:

\[
\operatorname {S i m} _ {V} = \psi (V) \cdot \psi (V) ^ {\top} \in \mathbb {R} ^ {n \times n}, \tag {7}
\]

where \(V ~ \in ~ \mathbb { R } ^ { n \times d _ { k } }\) denotes the
projected value featuresfrom concatenated input \(F _ { i n }\) , and
\(\psi ( \cdot )\) denotes row-wise L2-normalized transformation. A
binary mask matrix is thengenerated by thresholding \(S i m _ { V }\)
and applied to the atten-√tion weights
\(W _ { a } = \mathrm { s o f t m a x } ( Q _ { h } K _ { h } ^ { \top } / \sqrt { \dot { d } _ { k } } )\)
, thereby guidingthe attention mechanism to focus on features from
candidateregions that exhibit high similarity to the current target.

\section{2.2 Post-Processing}\label{post-processing}

In DSA sequences, a distinct lesion visible in one frame of-ten
corresponds to identifiable lesion regions in adjacent po-sitions of
neighboring frames. We refer to this pattern asspatiotemporal
consistency. Although our targeted featureaggregation enhances
detections, it cannot perfectly guar-antee spatiotemporal consistency of
results. The reason lies

\pandocbounded{\includegraphics[keepaspectratio,alt={image}]{images/25a08863d7f622f35886367b735d79786549b6d56e07f2162a3bb754fd6d983d.jpg}}

Figure 3: Bi-SeqNMS bounding box optimization workflow.Rescores boxes
through forward and reverse bounding boxlinking, followed by NMS to
eliminate redundancies andproduce final predictions.

in the FAB's self-attention mechanism, which relies on co-sine
similarity between features to aggregate information.In some cases, a
correct detection in a challenging sce-nario, such as an early frame
with poor contrast, might yielda feature vector with low similarity to
those from clearerframes. This could cause the attention mask to
erroneouslysuppress this valid proposal, leading to a missed detectionin
that frame and thus breaking the consistency. There-fore, to further
refine the results, we designed a Bidirec-tional Sequence Non-Maximum
Suppression (Bi-SeqNMS)method, inspired by Sequence Non-Maximum
Suppression(Seq-NMS) (Han et al.~2016), as illustrated in Fig. 3

The method first establishes continuous detection box tra-jectories
through temporal association. If the intersection-over-union (IoU)
between two bounding boxes in adjacentframes exceeds a predefined
threshold, they are identifiedas part of the same target trajectory.
Subsequently, sequencerescoring is applied to each trajectory by fusing
detectionconfidence scores across all frames within the trajectory.
Fi-nally, cross-frame redundancy suppression removes overlap-ping boxes
exceeding a threshold within the track, retain-ing only the
highest-scoring boundary box. Building uponthe foundation of Seq-NMS,
our approach performs anal-ogous inference on reverse-direction DSA
sequences andmerges results from forward and reverse inferences.
Image-level NMS is then applied to eliminate duplicate final out-puts.
By incorporating reverse-direction linking, detectionresults for the
current frame can be refined based on infor-mation from subsequent
frames.

\section{3 Experiments}\label{experiments}

The DSA sequence dataset was obtained through collabora-tion with a
hospital and consists of intraoperative and post-operative DSA images.
As introduced in Table 1, this datasetincludes angiographic views of the
internal carotid artery(ICA), external carotid artery (ECA), and
vertebral artery(VA). All images were anonymized by removing
clinicalmetadata to ensure compliance with patient privacy regula-

tions. All DICOM-format data were acquired using SiemensAXIOM Artis
angiography system (Siemens Healthineers,Germany) and Philips Allura
UNIQ system (Philips, Nether-lands) during routine neurointerventional
procedures.

The dataset for this study was collected from 78 patients,comprising a
total of 128 DSA sequences. These sequencesinclude 106 2D-DSA and 22
3D-RA, all with a resolution of\(1 0 2 4 \times 1 0 2 4\) pixels. The
dataset was prepared under the guid-ance of neurosurgeons. For
aneurysm-positive sequences,we retained only arterial phase frames that
clearly displaythe lesion, while also including sequences without IAs
asnegative samples. Consequently, the selected 2D-DSA se-quences range
from 10 to 16 frames, and 3D-RA sequencesspan 50 to 68 frames. To use
multi-angle vascular morphol-ogy, 22 original 3D-RA sequences were
processed into 25sub-sequences. To prevent data leakage, sub-sequences
fromthe same source were separated by an imaging angle greaterthan
\(6 0 ^ { \circ }\) , ensuring distinct vascular views. Combined withthe
106 2D-DSA sequences, this resulted in a final datasetof 131 sequences
used in our experiments. Bounding boxannotations for all IAs were
performed by a junior radiol-ogist and subsequently reviewed and refined
by two neuro-surgeons, each with over 5 years of experience in DSA
in-terpretation. All annotated data were then converted into theYOLO
dataset format. To standardize input dimensions, allsequences were
temporally resampled to 10 frames for net-work training. The entire
dataset was then evaluated using afive-fold cross-validation scheme to
ensure robust results.

\section{3.1 Implementation Details}\label{implementation-details}

The proposed method is implemented using Python 3.8.6and the PyTorch
2.0.1 framework, with model training con-ducted on an NVIDIA RTX4090 GPU
(24GB). The networkarchitecture utilizes YOLOv8-m as the backbone, with
boththe backbone and feature pyramid network (FPN) compo-nents
initialized using COCO pretrained weights from Ultr-alytics. We adopt a
stochastic gradient descent (SGD) opti-mizer with a momentum coefficient
of 0.937 and a weightdecay of 0.0005, combined with a linear learning
rate decaystrategy.

During the training process, the input image size wasmaintained at
\(1 0 2 4 \times 1 0 2 4\) pixels. We employed data aug-mentation
methods, including mosaic, random flips, ran-dom contrast adjustments
and random scaling, with the mo-saic augmentation being disabled for the
final 10 epochs.The training process consists of two phases. During
thefirst phase, the backbone network is trained on our DSAdataset to
establish baseline lesion detection capabilities.This phase employs a
maximum of 300 training epochs witha batch size of 8, starting with a
base learning rate of 0.01accompanied by a 10-epoch warm-up period. In
the sub-sequent phase, we introduce a feature aggregation modulewhile
maintaining the weights inherited from the first-phasemodel. This stage
is configured for 80 training epochs, pro-cesses sequence-consistent
images in unified batches (batchsize \(= 1\) ), and implements an
initial learning rate of 0.001with a 5-epoch warm-up period.

Principal hyperparameter settings were as follows. In theFAB, the cosine
similarity threshold for generating the at-

tention mask was set to 0.9. The general NMS thresholdfor eliminating
overlapping proposals was 0.75. For the Bi-SeqNMS during inference, the
IoU threshold for trajectorylinking and the final NMS suppression
threshold were set to0.20 and 0.95, respectively.

\section{3.2 Evaluation Metrics}\label{evaluation-metrics}

To evaluate model performance, we employed both local-ization and
classification metrics. For IA localization, weutilized COCO metrics
(Lin et al.~2014): Average Precision(AP), calculated over Intersection
over Union (IoU) thresh-olds ranging from 0.50 to 0.95; AP50 (AP at
\$\mathrm { I o U } \{ = \} 0 . 5 0 \$ ); andAP75 (AP at
\(\mathrm { I o U } { = } 0 . 7 5\) . For image-level classification,
animage was considered positive if it contained at least oneIA
detection, with its confidence score determined by thehighest score
among all detected boxes within that image.Images without an IA were
classified as the negative class.The primary classification metric was
the Area Under theReceiver Operating Characteristic Curve (AUC-ROC).
Wealso reported sensitivity (Sens) and specificity (Spec), de-fined as:
Sensitivity \(( \mathrm { S e n s } ) = \mathrm { T P } ~ /\)
\(( \mathrm { T P } + \mathrm { F N } )\) );
Specificity\(\mathrm { ( S p e c ) = T N / \left( T N + F P \right) }\)
\(( \mathrm { T N } + \mathrm { F P } )\) , where TP, FP, TN, and FN
repre-sent true positives, false positives, true negatives, and
falsenegatives, respectively. The optimal threshold for calculat-ing
these metrics was determined by maximizing Youden'sJ statistic. All
reported performance metrics are based on afive-fold cross-validation at
the image level.

\section{3.3 Comparative Experiments}\label{comparative-experiments}

In this section, we evaluate the proposed method using five-fold
cross-validation on the dataset. To validate the effective-ness of our
approach, we conduct quantitative comparisonswith advanced object
detection methods including Faster R-CNN (Ren et al.~2016), Cascade
R-CNN (Cai and Vascon-celos 2019), Swin (Liu et al.~2021), RetinaNet
(Ross andDollar 2017), FCOS (Tian et al.~2019), YOLOv8 (Jocher, ´Qiu,
and Chaurasia 2023) and DSA sequence-specific de-tection methods (Su et
al.~2022).

The quantitative comparison results, as shown in Ta-ble 2, demonstrate
that our method achieves superior per-formance across classification and
localization metrics com-pared to existing approaches. Among
single-frame detec-tors, YOLOv8 demonstrates superior localization
perfor-mance, surpassing the second-ranked Swin Transformer
by\(2 . 0 9 \%\) , \(2 9 . 6 0 \%\) , and \(5 2 . 6 5 \%\) in AP50, AP,
and AP75 met-rics respectively. However, Swin Transformer shows
betterperformance in classification metrics. After evaluating
lo-calization accuracy, classification performance, and modelparameters,
we have selected YOLOv8 as the baseline dueto its balance between these
factors. We further enhancedit with targeted improvements. By
integrating our proposedmodules, we achieved performance enhancements of
\(7 . 8 7 \%\) ,\(6 . 7 3 \%\) , \(5 . 4 2 \%\) , and \(9 . 7 2 \%\) in
AP50, AP, AP75, and AUCmetrics respectively over the original baseline,
surpassingSwin Transformer in classification performance. In compar-ison
to Su et al.~(2022)'s Conv-LSTM approach for wholefeature map temporal
feature extraction (with their featureextractor replaced by YOLOv8 for
fair comparison), our re-sults demonstrates the superiority of our
proposed feature

Table 2: Quantitative comparison of IA detection methods: Evaluation
metrics for localization and classification withmean±standard deviation
from five-fold cross-validation.

{\def\LTcaptype{none} % do not increment counter
\begin{longtable}[]{@{}lllllll@{}}
\toprule\noalign{}
\endhead
\bottomrule\noalign{}
\endlastfoot
\multirow{2}{*}{Methods} & \multicolumn{3}{l}{%
Localization Metrics} & \multicolumn{3}{l@{}}{%
Classification Metrics} \\
& AP50 & AP & AP75 & AUC & Spec & Sens \\
Faster R-CNN & 0.538 ± 0.162 & 0.298 ± 0.093 & 0.298 ± 0.090 & 0.789 ±
0.120 & 0.740 ± 0.101 & 0.763 ± 0.119 \\
Cascade R-CNN & 0.572 ± 0.156 & 0.329 ± 0.105 & 0.359 ± 0.142 & 0.796 ±
0.110 & 0.808 ± 0.046 & 0.702 ± 0.181 \\
Swin & 0.622 ± 0.144 & 0.321 ± 0.073 & 0.302 ± 0.071 & 0.847 ± 0.070 &
0.880 ± 0.074 & 0.719 ± 0.154 \\
RetinaNet & 0.551 ± 0.110 & 0.325 ± 0.065 & 0.348 ± 0.086 & 0.818 ±
0.076 & 0.880 ± 0.076 & 0.666 ± 0.159 \\
FCOS & 0.576 ± 0.123 & 0.312 ± 0.081 & 0.313 ± 0.091 & 0.838 ± 0.075 &
0.868 ± 0.076 & 0.731 ± 0.118 \\
YOLOv8 & 0.635 ± 0.137 & 0.416 ± 0.121 & 0.461 ± 0.149 & 0.813 ± 0.059 &
0.932 ± 0.103 & 0.618 ± 0.059 \\
Su et al. & 0.645 ± 0.113 & 0.393 ± 0.093 & 0.420 ± 0.125 & 0.834 ±
0.084 & 0.912 ± 0.117 & 0.726 ± 0.068 \\
Ours & 0.685 ± 0.132 & 0.444 ± 0.113 & 0.486 ± 0.136 & 0.892 ± 0.077 &
0.932 ± 0.083 & 0.806 ± 0.130 \\
\end{longtable}
}

\pandocbounded{\includegraphics[keepaspectratio,alt={image}]{images/847fd245b36d63550be94b5bfceaacb2d703a515b3c64b95d9b2c07d8ab92828.jpg}}

\pandocbounded{\includegraphics[keepaspectratio,alt={image}]{images/6a2544230a68c7c09e181ed780554cab1eb2dcc30b54966d000a590fa321104b.jpg}}

Figure 4: Visual comparison of IA detection results. In eachcase, the
top row presents the ground truth sequence. Thesubsequent rows display
detection visualizations from vari-ous models. Red arrows indicate
missed detections, and yel-low arrows indicate false detections.

aggregation enhancement, yielding more reliable and robustdetection
outcomes for clinical diagnostics.

To intuitively demonstrate the advantages of our method,Fig. 4
visualizes its performance in comparison to the base-

line (YOLOv8) and the primary sequence-based method (Suet al.). We
selected these representative methods for compar-ison to clearly
highlight the specific improvements in accu-racy and consistency
achieved by our method. In case (a),the baseline method did not detect
the IA in specific frameswhen faced with contrast flow or IA occlusion
by vessels.In contrast, our method provided more accurate and
consis-tent detection results. Case (b) exhibits an instance of
falsepositives from the baseline method, where a vessel overlap(yellow
arrow) was identified as an IA. Furthermore, case(b) also visualizes the
results of the method proposed bySu et al.~(2022). Although quantitative
metrics indicate thattheir method utilizes Conv-LSTM to effectively
correlatedetection results across frames, it also identified the
vesseloverlap as a false positive in this instance. The persistence
ofthis detection across frames suggests that its temporal corre-lation
may also apply to IA-like structures. In contrast, ourmethod effectively
addresses this issue. Through these twocases, we demonstrate that our
method not only maintainsstable detection under contrast flow and
occlusion caused byvessels, but also significantly reduces false
positives arisingfrom vessel overlaps that resemble IAs.

\section{3.4 Ablation Experiments}\label{ablation-experiments}

The proposed method integrates TFA into the base YOLOv8model, along with
a post-processing module named Bi-SeqNMS. To demonstrate the
contributions of each compo-nent in the proposed network, we conducted
ablation exper-iments using YOLOv8 as the baseline. The experimental
re-sults are presented in Table 3.

Ablation for TFA We first conducted an ablation study byintegrating the
Targeted Feature Aggregation (TFA) into theYOLOv8 baseline. As shown in
Row 2 of Table 3, incorpo-rating the TFA improved all detection metrics,
with AP50,AP, AP75, and AUC increasing by \(6 . 6 1 \%\) ,
\(5 . 5 3 \%\) , \(5 . 2 1 \%\) ,and \(8 . 3 4 \%\) , respectively. We
attribute this enhancement tothe FAB's capability to search for
supporting information fora target from other frames in the sequence.
This mechanismenables the model to re-evaluate the current frame based
ondetections in similar regions from other frames, particularlyin
instances of target degradation.

Additionally, we conducted controlled experiments to

Table 3: Ablation study of these components of the proposed method. The
results are average±standard deviation obtainedfrom the five-fold
cross-validation experiment.

{\def\LTcaptype{none} % do not increment counter
\begin{longtable}[]{@{}|l|l|l|l|l|l|l|l|l|@{}}
\toprule\noalign{}
\endhead
\bottomrule\noalign{}
\endlastfoot
\hline
\multicolumn{3}{@{}l}{%
Ablation Sets} & \multicolumn{3}{l}{%
Localization Metrics} & \multicolumn{3}{l@{}}{%
Classification Metrics} \\
\hline
Base & TFA & PPM & AP50 & AP & AP75 & AUC & Spec & Sens \\
\hline
✓ & & & 0.635 ± 0.137 & 0.416 ± 0.121 & 0.461 ± 0.149 & 0.813 ± 0.059 &
0.916 ± 0.121 & 0.686 ± 0.055 \\
\hline
✓ & ✓ & & 0.675 ± 0.134 & 0.439 ± 0.114 & 0.482 ± 0.139 & 0.887 ± 0.077
& 0.928 ± 0.058 & 0.784 ± 0.080 \\
\hline
✓ & & ✓ & 0.668 ± 0.122 & 0.444 ± 0.112 & 0.508 ± 0.132 & 0.831 ± 0.066
& 0.944 ± 0.071 & 0.716 ± 0.047 \\
\hline
✓ & ✓ & ✓ & 0.685 ± 0.132 & 0.444 ± 0.113 & 0.486 ± 0.136 & 0.892 ±
0.077 & 0.932 ± 0.083 & 0.806 ± 0.130 \\
\hline
\end{longtable}
}

Table 4: Ablation study on the impact of maximum numberof proposal
\(( P _ { m a x } )\) in feature selection process, with resultsaveraged
over five-fold cross-validation experiments.

{\def\LTcaptype{none} % do not increment counter
\begin{longtable}[]{@{}|l|l|l|l|l|@{}}
\toprule\noalign{}
\endhead
\bottomrule\noalign{}
\endlastfoot
\hline
Settings & AP50 & AP & AP75 & AUC \\
\hline
Pmax=3 & 0.660 & 0.429 & 0.478 & 0.866 \\
\hline
Pmax=5 & 0.670 & 0.431 & 0.471 & 0.887 \\
\hline
Pmax=10 & 0.675 & 0.439 & 0.482 & 0.887 \\
\hline
Pmax=15 & 0.667 & 0.428 & 0.462 & 0.883 \\
\hline
Pmax=20 & 0.665 & 0.429 & 0.469 & 0.882 \\
\hline
Pmax=30 & 0.670 & 0.430 & 0.467 & 0.879 \\
\hline
\end{longtable}
}

Table 5: Comparison between original and improved post-processing
modules, with results averaged over five-foldcross-validation
experiments. T, S, and B represent TFA,Seq-NMS, and Bi-SeqNMS,
respectively.

{\def\LTcaptype{none} % do not increment counter
\begin{longtable}[]{@{}|l|l|l|l|l|@{}}
\toprule\noalign{}
\endhead
\bottomrule\noalign{}
\endlastfoot
\hline
Settings & AP50 & AP & AP75 & AUC \\
\hline
Base & 0.635 & 0.416 & 0.461 & 0.813 \\
\hline
Base+S & 0.656 & 0.438 & 0.498 & 0.829 \\
\hline
Base+B & 0.668 & 0.444 & 0.508 & 0.831 \\
\hline
Base+T+S & 0.681 & 0.441 & 0.485 & 0.889 \\
\hline
Base+T+B & 0.685 & 0.444 & 0.486 & 0.892 \\
\hline
\end{longtable}
}

evaluate the impact of the maximum number of proposalsper frame
\(( P _ { m a x } )\) on detection performance. As shown inTable 4,
performance improved significantly across mostmetrics as
\(P _ { m a x }\) was increased from 3 to 10, reaching itsoptimum at
\(P _ { m a x } = 1 0\) . However, detection accuracy beganto decline as
\(P _ { m a x }\) exceeded 10. This phenomenon may beexplained by the
hypothesis that performance peaks once theselected proposals
sufficiently cover all targets detectable bythe baseline network. As
\(P _ { m a x }\) increases further, the inter-ference from redundant
background features on the aggre-gation process becomes more pronounced.

Ablation for PPM In this section, we evaluate the im-pact of the
post-processing module (PPM) on detection per-formance. We integrated
Bi-SeqNMS into two frameworks:one without TFA and one with TFA, while
maintaining con-sistent experimental settings. The quantitative results
in Ta-ble 3 demonstrate that the integration of Bi-SeqNMS sig-nificantly
enhances detection performance compared to theframework without PPM.
When both TFA and Bi-SeqNMS

were used, the AP50, AP, AP75, and AUC of the test resultswere 0.685,
0.444, 0.486, and 0.892, respectively. Comparedto the framework without
PPM, each metric improved by\(1 . 4 8 \%\) , \(1 . 1 4 \%\) ,
\(0 . 8 3 \%\) , and \(0 . 5 6 \%\) , respectively.

Furthermore, we experimentally evaluate the effective-ness of the
enhancements made to Seq-NMS. By applyingSeq-NMS and Bi-SeqNMS
post-processing to detection re-sults from the same trained model, Table
5 demonstratesthat both methods enhance detection metrics compared tothe
absence of post-processing, whether TFA is utilized ornot. The optimal
performance was achieved when the cur-rent framework employed Bi-SeqNMS,
yielding AP50, AP,AP75, and AUC values of 0.685, 0.444, 0.486, and
0.892,respectively. AP50, AP, AP75, and AUC were further im-proved by
\(0 . 5 8 \%\) , \(0 . 6 8 \%\) , \(0 . 2 1 \%\) , and \(0 . 3 4 \%\) ,
respectivelycompared to the use of Seq-NMS. This performance advan-tage
may stem from Bi-SeqNMS's optimization of early-sequence detection. In
early DSA sequences where IA ves-sels are not fully filled with contrast
agent, detection scorestend to be lower. Since Seq-NMS relies on
preceding framesto rescore current frames, it limits improvements in
earlyframes. Bi-SeqNMS compensates through reverse linkingoperations;
therefore, this bidirectional approach enablesbetter optimization than
Seq-NMS.

\section{4 Conclusion}\label{conclusion}

In this study, we present a deep learning framework designedto enhance
the spatiotemporal consistency of IA detection,applicable to both 2D-DSA
and 3D-RA images. Our workaddresses two principal challenges: (1) the
spatial scale dis-parity between small lesion areas and vast background
fea-tures, and (2) the frequent lack of spatiotemporal consis-tency in
detection results. The TFA we proposed is essen-tial for achieving this
consistency enhancement. First, its in-tegrated feature selection
mechanism effectively suppressesbackground noise by dynamically
filtering for high-qualityproposals. Subsequently, FAB utilizes these
targeted fea-tures from across the sequence to build a temporally
sta-ble and robust representation for underlying lesions. Finally,our
proposed Bi-SeqNMS module further refines this con-sistency through
bidirectional temporal linking. Extensiveexperimental evaluations
demonstrate that our frameworkachieves superior performance in IA
detection tasks com-pared to existing methods when applied to DSA
sequences.These results suggest the potential clinical utility of our
sys-tem as a decision support tool for IA screening applications.

\section{References}\label{references}

Cai, Z.; and Vasconcelos, N. 2019. Cascade R-CNN: Highquality object
detection and instance segmentation. IEEEtransactions on pattern
analysis and machine intelligence,43(5): 1483--1498.

Duan, H.; Huang, Y.; Liu, L.; Dai, H.; Chen, L.; and Zhou,L. 2019.
Automatic detection on intracranial aneurysm fromdigital subtraction
angiography with cascade convolutionalneural networks. Biomedical
engineering online, 18: 1--18.

Faron, A.; Sichtermann, T.; Teichert, N.; Luetkens, J. A.;Keulers, A.;
Nikoubashman, O.; Freiherr, J.; Mpotsaris, A.;and Wiesmann, M. 2020.
Performance of a deep-learningneural network to detect intracranial
aneurysms from 3DTOF-MRA compared to human readers. Clinical
neurora-diology, 30: 591--598.

Feng, C.; Zhong, Y.; Gao, Y.; Scott, M. R.; and Huang, W.2021. Tood:
Task-aligned one-stage object detection. In2021 IEEE/CVF International
Conference on Computer Vi-sion (ICCV), 3490--3499. IEEE Computer
Society.

Han, W.; Khorrami, P.; Paine, T. L.; Ramachandran, P.;Babaeizadeh, M.;
Shi, H.; Li, J.; Yan, S.; and Huang, T. S.2016. Seq-nms for video object
detection. arXiv preprintarXiv:1602.08465.

Hiratsuka, Y.; Miki, H.; Kiriyama, I.; Kikuchi, K.; Taka-hashi, S.;
Matsubara, I.; Sadamoto, K.; and Mochizuki, T.2008. Diagnosis of
unruptured intracranial aneurysms: 3TMR angiography versus 64-channel
multi-detector row CTangiography. Magnetic Resonance in Medical
Sciences,7(4): 169--178.

Jocher, G.; Qiu, J.; and Chaurasia, A. 2023. YOLOv8 by Ul-tralytics.
https://github.com/ultralytics/ultralytics. Version8.1.27.

Li, M.-H.; Chen, S.-W.; Li, Y.-D.; Chen, Y.-C.; Cheng, Y.-S.; Hu, D.-J.;
Tan, H.-Q.; Wu, Q.; Wang, W.; Sun, Z.-K.;et al.~2013. Prevalence of
unruptured cerebral aneurysms inChinese adults aged 35 to 75 years: a
cross-sectional study.Annals of internal medicine, 159(8): 514--521.

Lin, T.-Y.; Maire, M.; Belongie, S.; Hays, J.; Perona, P.; Ra-manan, D.;
Dollar, P.; and Zitnick, C. L. 2014. Microsoft ´coco: Common objects in
context. In Computer vision--ECCV 2014: 13th European conference,
zurich, Switzer-land, September 6-12, 2014, proceedings, part v 13,
740--755. Springer.

Liu, Z.; Lin, Y.; Cao, Y.; Hu, H.; Wei, Y.; Zhang, Z.; Lin,S.; and Guo,
B. 2021. Swin transformer: Hierarchical vi-sion transformer using
shifted windows. In Proceedings ofthe IEEE/CVF international conference
on computer vision,10012--10022.

Mo, Y.; Chen, Y.; Hu, Y.; Xu, J.; Wang, H.; Dai, L.; and Liu,J. 2023.
Focusing Intracranial Aneurysm Lesion Segmenta-tion by Graph Mask2Former
with Local Refinement in DSAImages. In 2023 IEEE International
Conference on Bioin-formatics and Biomedicine (BIBM), 899--903. IEEE.

Piotin, M.; Gailloud, P.; Bidaut, L.; Mandai, S.; Muster, M.;Moret, J.;
and Rufenacht, D. A. 2003. CT angiography, MR ¨angiography and
rotational digital subtraction angiography

for volumetric assessment of intracranial aneurysms. An ex-perimental
study. Neuroradiology, 45: 404--409.

Ren, S.; He, K.; Girshick, R.; and Sun, J. 2016. Faster R-CNN: Towards
real-time object detection with region pro-posal networks. IEEE
transactions on pattern analysis andmachine intelligence, 39(6):
1137--1149.

Rodriguez-Regent, C.; Edjlali-Goujon, M.; Trystram, D.; ´Boulouis, G.;
Hassen, W. B.; Godon-Hardy, S.; Nataf, F.;Machet, A.; Legrand, L.;
Ladoux, A.; et al.~2014. Non-invasive diagnosis of intracranial
aneurysms. Diagnosticand interventional imaging, 95(12): 1163--1174.

Ropper, A. H.; and Zervas, N. T. 1984. Outcome 1 yearafter SAH from
cerebral aneurysm: management morbidity,mortality, and functional status
in 112 consecutive good-riskpatients. Journal of neurosurgery, 60(5):
909--915.

Ross, T.-Y.; and Dollar, G. 2017. Focal loss for dense ob- ´ject
detection. In proceedings of the IEEE conference oncomputer vision and
pattern recognition, 2980--2988.

Shen, D.; Wu, G.; and Suk, H.-I. 2017. Deep learning inmedical image
analysis. Annual review of biomedical engi-neering, 19(1): 221--248.

Shi, Y.; Wang, N.; and Guo, X. 2023. YOLOV: Makingstill image object
detectors great at video object detection.In Proceedings of the AAAI
conference on artificial intelli-gence, volume 37, 2254--2262.

Shi, Z.; Miao, C.; Schoepf, U. J.; Savage, R. H.; Dargis,D. M.; Pan, C.;
Chai, X.; Li, X. L.; Xia, S.; Zhang, X.; et al.2020. A clinically
applicable deep-learning model for de-tecting intracranial aneurysm in
computed tomography an-giography images. Nature communications, 11(1):
6090.

Su, R.; van der Sluijs, M.; Cornelissen, S. A.; Lycklama, G.;Hofmeijer,
J.; Majoie, C. B.; van Doormaal, P. J.; Van Es,A. C.; Ruijters, D.;
Niessen, W. J.; et al.~2022. Spatio-temporal deep learning for automatic
detection of intracra-nial vessel perforation in digital subtraction
angiographyduring endovascular thrombectomy. Medical image anal-ysis,
77: 102377.

Tian, Z.; Shen, C.; Chen, H.; and He, T. 2019. Fcos: Fullyconvolutional
one-stage object detection. In Proceedings ofthe IEEE/CVF international
conference on computer vision,9627--9636.

Vaswani, A.; Shazeer, N.; Parmar, N.; Uszkoreit, J.; Jones,L.; Gomez, A.
N.; Kaiser, Ł.; and Polosukhin, I. 2017. At-tention is all you need.
Advances in neural information pro-cessing systems, 30.

Wu, H.; Chen, Y.; Wang, N.; and Zhang, Z. 2019. Sequencelevel semantics
aggregation for video object detection. InProceedings of the IEEE/CVF
international conference oncomputer vision, 9217--9225.

Xu, W.; Yang, H.; Tan, T.; Liu, W.; Gao, F.; Deng, Y.; Pan,X.; and Su,
R. 2023. Improved YOLOX Framework forAutomatic Large Intracranial Artery
Stenosis Detection inDigital Subtraction Angiography Images. In 2023
IEEE In-ternational Conference on Bioinformatics and Biomedicine(BIBM),
2326--2331. IEEE Computer Society.

Yoon, N. K.; McNally, S.; Taussky, P.; and Park, M. S. 2016.Imaging of
cerebral aneurysms: a clinical perspective. Neu-rovascular Imaging,
2(1): 6.

Zeng, Y.; Liu, X.; Xiao, N.; Li, Y.; Jiang, Y.; Feng, J.; andGuo, S.
2019. Automatic diagnosis based on spatial infor-mation fusion feature
for intracranial aneurysm. IEEE trans-actions on medical imaging, 39(5):
1448--1458.

Zhang, J.; Xie, Q.; Mou, L.; Zhang, D.; Chen, D.; Shan,C.; Zhao, Y.; Su,
R.; and Guo, M. 2025. DSCA: A Digi-tal Subtraction Angiography Sequence
Dataset and Spatio-Temporal Model for Cerebral Artery Segmentation.
IEEETransactions on Medical Imaging.

Zhang, J.; Zhao, Y.; Liu, X.; Jiang, J.; and Li, Y. 2023.FSTIF-UNet: a
deep learning-based method towards au-tomatic segmentation of
intracranial aneurysms in un-reconstructed 3D-RA. IEEE Journal of
Biomedical andHealth Informatics, 27(8): 4028--4039.

Zheng, Z.; Wang, P.; Ren, D.; Liu, W.; Ye, R.; Hu, Q.; andZuo, W. 2021.
Enhancing geometric factors in model learn-ing and inference for object
detection and instance segmen-tation. IEEE transactions on cybernetics,
52(8): 8574--8586.

\section{Reproducibility Checklist}\label{reproducibility-checklist}

\section{1. General Paper Structure}\label{general-paper-structure}

1.1. Includes a conceptual outline and/or pseudocode de-scription of AI
methods introduced (yes/partial/no/NA)yes

1.2. Clearly delineates statements that are opinions, hypoth-esis, and
speculation from objective facts and results(yes/no) yes

1.3. Provides well-marked pedagogical references for less-familiar
readers to gain background necessary to repli-cate the paper (yes/no)
yes

\section{2. Theoretical Contributions}\label{theoretical-contributions}

2.1. Does this paper make theoretical contributions?(yes/no) no

If yes, please address the following points:

2.2. All assumptions and restrictions are stated clearlyand formally
(yes/partial/no) NA

2.3. All novel claims are stated formally (e.g., in theoremstatements)
(yes/partial/no) NA

2.4. Proofs of all novel claims are included (yes/par-tial/no) NA

2.5. Proof sketches or intuitions are given for complexand/or novel
results (yes/partial/no) NA

2.6. Appropriate citations to theoretical tools used aregiven
(yes/partial/no) NA

2.7. All theoretical claims are demonstrated empiricallyto hold
(yes/partial/no/NA) NA

2.8. All experimental code used to eliminate or disproveclaims is
included (yes/no/NA) NA

\section{3. Dataset Usage}\label{dataset-usage}

3.1. Does this paper rely on one or more datasets? (yes/no)yes

If yes, please address the following points:

3.2. A motivation is given for why the experimentsare conducted on the
selected datasets (yes/par-tial/no/NA) yes

3.3. All novel datasets introduced in this paper are in-cluded in a data
appendix (yes/partial/no/NA) yes

3.4. All novel datasets introduced in this paper will bemade publicly
available upon publication of the pa-per with a license that allows free
usage for researchpurposes (yes/partial/no/NA) no

3.5. All datasets drawn from the existing literature (po-tentially
including authors' own previously pub-lished work) are accompanied by
appropriate cita-tions (yes/no/NA) NA

3.6. All datasets drawn from the existing literature(potentially
including authors' own previouslypublished work) are publicly available
(yes/par-tial/no/NA) NA

3.7. All datasets that are not publicly available are de-scribed in
detail, with explanation why publiclyavailable alternatives are not
scientifically satisficing(yes/partial/no/NA) yes

\section{4. Computational Experiments}\label{computational-experiments}

4.1. Does this paper include computational experiments?(yes/no) yes

If yes, please address the following points:

4.2. This paper states the number and range of valuestried per (hyper-)
parameter during development ofthe paper, along with the criterion used
for selectingthe final parameter setting (yes/partial/no/NA) par-tial

4.3. Any code required for pre-processing data is in-cluded in the
appendix (yes/partial/no) yes

4.4. All source code required for conducting and analyz-ing the
experiments is included in a code appendix(yes/partial/no) yes

4.5. All source code required for conducting and ana-lyzing the
experiments will be made publicly avail-able upon publication of the
paper with a license

that allows free usage for research purposes (yes/-partial/no) no

4.6. All source code implementing new methods havecomments detailing the
implementation, with refer-ences to the paper where each step comes from
(yes/-partial/no) yes

4.7. If an algorithm depends on randomness, then themethod used for
setting seeds is described in a waysufficient to allow replication of
results (yes/par-tial/no/NA) yes

4.8. This paper specifies the computing infrastructureused for running
experiments (hardware and soft-ware), including GPU/CPU models; amount
ofmemory; operating system; names and versions ofrelevant software
libraries and frameworks (yes/par-tial/no) yes

4.9. This paper formally describes evaluation metricsused and explains
the motivation for choosing thesemetrics (yes/partial/no) yes

4.10. This paper states the number of algorithm runs usedto compute each
reported result (yes/no) yes

4.11. Analysis of experiments goes beyond single-dimensional summaries
of performance (e.g., aver-age; median) to include measures of
variation, con-fidence, or other distributional information (yes/no)yes

4.12. The significance of any improvement or decrease inperformance is
judged using appropriate statisticaltests (e.g., Wilcoxon signed-rank)
(yes/partial/no)yes

4.13. This paper lists all final (hyper-)parameters usedfor each
model/algorithm in the paper's experiments(yes/partial/no/NA) partial

\end{document}
